\scanpage{27}

\margin{1} $\underline{\overline{\mathrm{XVIII}}}$ $1.2$ \mspage{8}

et si $G^n$ est l'image réciproque de $ {}_nA$, le groupe $H^1(X,G)$
est réunion des groupes $H^1(X,G^n)$. Les groupes $G^n$ sont
affins (non nécessairement lisse) et, sous l'hypothèse faite,
l'assertion résulte de $1.2.\deleted{\unsure{2|1|7}.}8$.

c. Pour $S$ strictement hensélien et $G$ lisse, on a $H^1(S,G)=0$.
il nous faut donc vérifier que $H^1(X,G)=0$. Supposons que
ce soit déjà établi après le changement de base $i:S_0\longrightarrow S$, où
$S_0$ est défini par un idéal de carré nul \inserted{et de longueur $1$} $I$ \deleted{\unsure{du|de}} de $\unsure{O|\mathcal O}_S$. Les
$G$-torseurs sur $X$ sont alors des déformations du $G$-torseur
trivial sur $X_0=X\times_S S_0$. Si $L=e^*\Omega^1_{G/S}$ ($e$ section neutre)
ces déformations sont classifiées par
\[
\deleted{H^1(S_0}
\]
\[
H^1(X_{\deleted{\unsure{0|s}}},f^*\underline{\mathrm{Hom}}(L,I))=0\ ,
\]
et ceci prouve $1.2.\deleted{\unsure{1|2|7}.}9$.
% archive sheet number: 26

